\section{Diskussion}
Die Experimente belegen einzelne Eigenschaften der entwickelten Software und Modelle, erlauben aber noch keine Aussage über ein vollständiges autonomes Sortiersystem. Insbesondere sind Validierungsmetriken auf kuratierten Datensätzen, Messungen auf einer Intel-CPU und Prognosen für die Zielhardware getrennt zu betrachten.

\subsection{Bewertung der ursprünglichen Hypothesen}
Die zu Projektbeginn formulierte Transfer-Learning-Hypothese H1 wird durch die historischen Werte allenfalls vorläufig gestützt: MobileNetV2 mit Fine-Tuning erreichte \SI{87.42}{\percent}, das Scratch-CNN \SI{61.00}{\percent}. EXP-001 entstand jedoch mit einer kleineren Datengrundlage. Der Abstand von 26,42 Prozentpunkten ist deshalb kein kontrollierter Architektureffekt und bestätigt H1 nicht kausal \cite{yosinski2014transferable}.

H2 ist nur hinsichtlich Genauigkeit und Dateigröße des Stage-A-Klassifikators auswertbar. Das Keras-Quellmodell umfasst \SI{23.48}{\mebi\byte}. Der Export in das TFLite-Format umfasst in FP32 \SI{8.49}{\mebi\byte} und in INT8 \SI{2.61}{\mebi\byte}. Gegenüber FP32-TFLite sank die Größe durch INT8 somit um rund \SI{69}{\percent}; die protokollierte Validierungsgenauigkeit sank um 0,07 Prozentpunkte. Die früher notierten Laufzeiten stammen nicht aus einer ausreichend dokumentierten, formatgleichen Messung und werden daher nicht zur Bestätigung einer Latenzhypothese herangezogen. Der Effekt auf den YOLO-Detektor und einen Raspberry Pi bleibt offen.

H3 ist vorläufig gestützt: Die Software kann mit YOLO11n mehrere Bounding Boxes in einem Bild ausgeben. EXP-014 erreichte auf seinem Validierungssplit \SI{60.7}{\percent} mAP50, für \textit{Trash} jedoch nur \SI{26.9}{\percent}. Damit ist die technische Mehrfachdetektion gezeigt, nicht aber eine robuste Lokalisierung in unbekannten Einsatzumgebungen. E1 ist nicht umgesetzt: Der aktuelle Inferenzpfad enthält keinen EMA-Filter. Eine Wirkung zeitlicher Glättung auf mechanische Fehltrigger oder einen Sortiervorgang wurde daher nicht gemessen.

\subsection{Gültigkeit der Auswertung}
EXP-005 ist kein positiver Leistungsnachweis. Die automatisch erzeugten Bounding Boxes enthielten systematisch Tischkanten; die hohe interne mAP50 von \SI{82.3}{\percent} entstand durch Data Leakage und Hintergrund-Overfitting. Das Versagen auf anderen Hintergründen zeigt, dass das Modell keine übertragbare Objektrepräsentation gelernt hatte.

Auch der als \enquote{Field Benchmark} bezeichnete Vergleich ist in seiner jetzigen Form \textbf{kein gültiger Feldtest}: Er verwendet das \texttt{mira\_v2}-Validierungsset und damit keine unabhängig im späteren Einsatz erhobenen Feldaufnahmen. Zudem ist die Auswertung mit einer einheitlichen Konfidenzschwelle für unterschiedlich kalibrierte FP32- und INT8-Modelle nicht hinreichend kontrolliert. Die dort berichteten F1-Werte, einschließlich \SI{85.8}{\percent} für EXP-014 und des scheinbaren INT8-Rückgangs, sind daher vorläufige Pipeline-Ausgaben. Sie dürfen weder als Generalisierungsnachweis noch als Beleg praktischer Sortierqualität interpretiert werden. Erforderlich sind ein vorab festgelegtes, unabhängiges Testset, identische Zuordnungsregeln, separat kalibrierte Schwellen und eine klassenweise Fehleranalyse.

Es liegen historische Laufzeitangaben von einer Intel-Core-i7-CPU und einer NVIDIA-T4-Cloud-GPU vor; mangels eines einheitlich dokumentierten Messprotokolls eignen sie sich jedoch nicht für den aktuellen FP32-/INT8-Vergleich. Die für Raspberry Pi genannten \SI{90}{\milli\second} bis \SI{120}{\milli\second} sind Projektionen, keine Messwerte. Weder Raspberry Pi noch Kamera-Roboter-Gesamtsystem wurden vermessen; Aussagen zu Bildrate, Greifgenauigkeit, Durchsatz, Energiebedarf oder Fehlwürfen sind deshalb derzeit nicht möglich.

\subsection{Klassen, Materialbegriff und Datenbias}
Die fünf Zielklassen sind visuelle Entsorgungskategorien, keine verlässliche stoffliche Analyse. Eine RGB-Kamera erfasst Form, Farbe, Textur und Reflexion, aber keine chemische Zusammensetzung. Weißes zerknülltes Papier kann deshalb wie Kunststofffolie erscheinen; äußerlich ähnliche Behälter können aus verschiedenen Werkstoffen bestehen. Mehrschichtige Verbunde, beschichtetes Papier, Getränkekartons und verschmutzte Verpackungen besitzen zudem keine eindeutige Zuordnung zu genau einer Materialklasse. \mira erkennt somit gelernte visuelle Kategorien und kann echte Materialidentität ohne zusätzliche Sensorik nicht garantieren.

Das Remapping der Quelldatensätze verschärft diese Einschränkung. Spezifische Objekt- und Materialtaxonomien wurden auf die fünf Zielklassen verdichtet.

Dabei gehen Unterschiede zwischen Objektart, Material und lokalem Entsorgungsweg verloren; Verbunde werden durch eine Projektentscheidung einer einzigen Klasse zugeordnet. Stage A enthält nur vier Klassen und muss unbekannten Restmüll zwangsläufig falsch zuordnen. In Stage B umfasst \textit{Trash} visuell sehr verschiedene Gegenstände. Auch im aktuellen Datensatz können Quellenabhängigkeit, kanonische Objektansichten, uneinheitliche Annotationen und inhaltlich ähnliche Aufnahmen über Splitgrenzen hinweg die Metriken verzerren \cite{ben2010domain}. Die historischen Läufe EXP-013 bis EXP-017 sprechen dafür, dass die Datenzusammensetzung die klassenweisen Ergebnisse stark beeinflusst; wegen unterschiedlicher Datensätze und Validierungssplits bilden sie aber keine kontrollierte Ablation.

\subsection{Stand der laufenden Arbeiten}
Der aktuelle balancierte Datensatz ist erstellt, das darauf vorgesehene YOLO11n-Training und die formatgleiche FP32-/INT8-Auswertung sind jedoch noch nicht abgeschlossen. Teacher-, Baseline- und Distillationsläufe befinden sich ebenfalls nur in Planung oder Bearbeitung. Erwartete mAP-Werte, Modellgrößen und Raspberry-Pi-Leistungen sind Projektionen und werden nicht als Ergebnisse verwendet. Erst nach abgeschlossenem Training und Evaluation auf dem eingefrorenen Testset ist ein belastbarer Vergleich möglich.
